% =====================================================================
%  IEL04 - Circuitos Eléctricos I
%  Semana 13: Cálculo y medición de parámetros en circuitos RLC en serie, en paralelo y mixtos: filtros resonantes, frecuencia de resonancia, ancho de banda y nomenclatura
%  Institución Universitaria de Barranquilla (IUB)
%  Generada por src/presentaciones.py (diseño IUB 5). Compilador: pdfLaTeX (dos pasadas)
% =====================================================================
\documentclass[aspectratio=169,11pt,xcolor={table}]{beamer}

% ---------------------------------------------------------------------
%  Codificación, idioma y tipografía
% ---------------------------------------------------------------------
\usepackage[utf8]{inputenc}
\usepackage[T1]{fontenc}
\usepackage[spanish,es-nodecimaldot,es-noshorthands]{babel}
\usepackage{avant}                       % Sans geométrica (emula Avenir Next)
\renewcommand{\familydefault}{\sfdefault}
\renewcommand{\ttdefault}{pcr}           % Monoespaciada para código
\usepackage{textcomp}

% ---------------------------------------------------------------------
%  Paquetes
% ---------------------------------------------------------------------
\usepackage{amsmath,amssymb,mathtools}
\usepackage{siunitx}
\sisetup{output-decimal-marker={.}, per-mode=symbol}
\usepackage{booktabs,tabularx,multirow,array}
\usepackage{adjustbox}
\usepackage{tikz}
\usetikzlibrary{arrows.meta,positioning,calc,shapes.geometric,shapes.misc,
  decorations.pathreplacing,fit,backgrounds,babel}
\usepackage[siunitx,RPvoltages]{circuitikz}
\usepackage{pgfplots}
\pgfplotsset{compat=1.18}
\usepackage[most]{tcolorbox}
\usepackage{listings}

% ---------------------------------------------------------------------
%  Paleta institucional IUB (Manual de Identidad Visual 2025)
% ---------------------------------------------------------------------
\definecolor{iubwhite}{HTML}{FFFFFF}
\definecolor{iubnavy}{HTML}{121023}
\definecolor{iubyellow}{HTML}{FFDF2D}
\definecolor{iubblue}{HTML}{00ADE7}
\definecolor{iuborange}{HTML}{E39037}
\definecolor{iubgreen}{HTML}{3B9C6D}

% ---------------------------------------------------------------------
%  Configuración de Beamer
% ---------------------------------------------------------------------
\setbeamertemplate{navigation symbols}{}
\setbeamersize{text margin left=0.6cm, text margin right=0.6cm}
\setbeamercolor{normal text}{fg=iubnavy, bg=iubwhite}
\setbeamercolor{background canvas}{bg=iubwhite}
\setbeamercolor{structure}{fg=iubnavy}
\setbeamercolor{alerted text}{fg=iuborange}
\setbeamertemplate{itemize item}{\textcolor{iubblue}{\small$\blacktriangleright$}}
\setbeamertemplate{itemize subitem}{\textcolor{iuborange}{\scriptsize$\bullet$}}
\setbeamertemplate{enumerate item}{\textcolor{iubblue}{\bfseries\insertenumlabel.}}
\setbeamertemplate{frametitle}{}         % el título vive en el headline

% Parches: el headline se pre-mide antes del primer frame
\makeatletter
\providecommand{\insertframetitle}{}
\providecommand{\insertframesubtitle}{}
\makeatother

% ---------- Encabezado ----------
\setbeamertemplate{headline}{%
  \begin{tikzpicture}
    \useasboundingbox (0,0) rectangle (\paperwidth,1.05cm);
    \fill[iubnavy] (0,0) rectangle (\paperwidth,1.05cm);
    \fill[iubyellow] (0,0) rectangle (\paperwidth,0.05cm);
    \node[anchor=west, text=iubyellow, font=\large\bfseries]
      at (0.6cm,0.55cm) {\strut\insertframetitle};
    \node[anchor=east, text=iubyellow, font=\footnotesize\bfseries]
      at ([xshift=-0.6cm]\paperwidth,0.55cm) {IUB};
  \end{tikzpicture}}

% ---------- Pie de página ----------
\setbeamertemplate{footline}{%
  \begin{tikzpicture}
    \useasboundingbox (0,0) rectangle (\paperwidth,0.55cm);
    \fill[iubnavy] (0,0) rectangle (\paperwidth,0.55cm);
    \node[anchor=west, text=iubyellow, font=\tiny]
      at (0.6cm,0.275cm) {IEL04 \textperiodcentered{} Circuitos Eléctricos I};
    \node[anchor=center, text=iubyellow, font=\tiny]
      at (0.66\paperwidth,0.275cm) {Ing. Sergio Martinez Castilla};
    \node[anchor=east, text=iubyellow, font=\tiny\bfseries]
      at ([xshift=-0.6cm]\paperwidth,0.275cm)
      {Semana 13 \quad \insertframenumber\,/\,\inserttotalframenumber};
  \end{tikzpicture}}

% ---------------------------------------------------------------------
%  Cajas tcolorbox (texto blanco sobre la paleta complementaria)
%  \begin{caja}[opciones]{color}{título} ... \end{caja}
%  \begin{cajaplana}[opciones]{color} ... \end{cajaplana}
%  \begin{cardB|cardO|cardG|cardN}[opciones]{título} ... \end{cardX}
% ---------------------------------------------------------------------
\tcbset{
  iubbase/.style={enhanced, arc=1.5mm, boxrule=0pt, coltext=white, coltitle=white,
    fonttitle=\bfseries\footnotesize, fontupper=\footnotesize,
    left=1.8mm, right=1.8mm, top=1mm, bottom=1.2mm,
    toptitle=0.6mm, bottomtitle=0.6mm, halign=left,
    before upper={\hyphenpenalty=5000\setbeamercolor{itemize item}{fg=white}%
                  \setbeamercolor{itemize/enumerate body}{fg=white}%
                  \setbeamercolor{itemize/enumerate subbody}{fg=white}%
                  \setbeamercolor{itemize subitem}{fg=white}%
                  \setbeamercolor{enumerate item}{fg=white}%
                  \setbeamertemplate{itemize item}{\textcolor{white}{\small$\blacktriangleright$}}}}
}
\newtcolorbox{caja}[3][]{iubbase, colback=#2, colframe=#2,
  colbacktitle=#2!78!iubnavy, title={#3}, #1}
\newtcolorbox{cajaplana}[2][]{iubbase, colback=#2, colframe=#2, #1}
\newtcolorbox{cardB}[2][]{iubbase, colback=iubblue,   colframe=iubblue,
  colbacktitle=iubblue!78!iubnavy,   title={#2}, #1}
\newtcolorbox{cardO}[2][]{iubbase, colback=iuborange, colframe=iuborange,
  colbacktitle=iuborange!78!iubnavy, title={#2}, #1}
\newtcolorbox{cardG}[2][]{iubbase, colback=iubgreen,  colframe=iubgreen,
  colbacktitle=iubgreen!78!iubnavy,  title={#2}, #1}
\newtcolorbox{cardN}[2][]{iubbase, colback=iubnavy,   colframe=iubnavy,
  colbacktitle=iubnavy, coltitle=iubyellow, title={#2}, #1}

% Celda de encabezado de tabla (texto amarillo sobre \rowcolor{iubnavy}) y código en línea
\newcommand{\hd}[1]{\textcolor{iubyellow}{\bfseries #1}}
\newcommand{\thc}[1]{\textcolor{iubyellow}{\textbf{#1}}}
\newcommand{\cod}[1]{\texttt{\textbf{#1}}}

% ---------------------------------------------------------------------
%  Código sobre fondo oscuro:
%  \begin{codebox}[listing options app={language=Python,firstnumber=1}]{título}
% ---------------------------------------------------------------------
\lstdefinestyle{iubcode}{
  basicstyle=\ttfamily\fontsize{7}{8.4}\selectfont\color{white},
  keywordstyle=\color{iubblue}\bfseries,
  commentstyle=\color{iubgreen!55!white}\itshape,
  stringstyle=\color{iuborange!80!white},
  numbers=left, numberstyle=\tiny\color{iubyellow}, numbersep=6pt,
  showstringspaces=false, breaklines=true, tabsize=4,
  columns=fullflexible, keepspaces=true, upquote=true}
\newtcblisting{codebox}[2][]{enhanced, listing only, colback=iubnavy,
  colframe=iubnavy, colbacktitle=iubnavy!85!white, coltitle=iubyellow,
  fonttitle=\bfseries\scriptsize, title={#2}, arc=1.5mm, boxrule=0pt,
  left=6mm, right=1mm, top=0.5mm, bottom=0.5mm, toptitle=0.5mm, bottomtitle=0.5mm,
  listing options={style=iubcode}, #1}

% ---------------------------------------------------------------------
%  Estilos TikZ
% ---------------------------------------------------------------------
\tikzset{
  bloque/.style={rectangle, rounded corners=2mm, fill=#1, text=white,
    font=\scriptsize\bfseries, align=center, minimum height=1.1cm,
    text width=1.8cm, inner sep=1.2mm},
  flecha/.style={-{Stealth[length=2.2mm]}, line width=1pt, draw=iubnavy},
  arr/.style={flecha},
  term/.style={rectangle, draw=#1, line width=1.2pt, fill=white,
    text=iubnavy, font=\tiny\bfseries, minimum width=1.05cm,
    minimum height=0.5cm, inner sep=1pt},
  nodo/.style={rectangle, rounded corners=1mm, fill=#1, text=white,
    font=\scriptsize\bfseries, minimum size=0.75cm, align=center},
  cable/.style={line width=1.4pt, draw=#1},
  box/.style={rectangle, rounded corners=1.5mm, draw=none, text=white,
    font=\scriptsize\bfseries, align=center, minimum height=0.8cm, inner sep=3pt},
  bB/.style={box, fill=iubblue}, bO/.style={box, fill=iuborange},
  bG/.style={box, fill=iubgreen}, bN/.style={box, fill=iubnavy},
  dec/.style={diamond, aspect=1.9, fill=iuborange, text=white,
    font=\scriptsize\bfseries, align=center, inner sep=1pt},
  tag/.style={fill=#1, text=white, font=\scriptsize\bfseries, rounded corners=1mm,
    minimum width=0.7cm, anchor=north west},
}

% ---------------------------------------------------------------------
%  Diapositiva separadora de bloque
%  #1 número  #2 título  #3 duración  #4 color
% ---------------------------------------------------------------------
\newcommand{\bloque}[4]{%
\section{#2}
\begin{frame}[plain]
\begin{tikzpicture}[remember picture, overlay]
  \fill[#4] (current page.north west) rectangle
        ([xshift=5.2cm]current page.south west);
  \node[anchor=north west, text=white, font=\bfseries\large]
    at ([xshift=0.8cm,yshift=-2.2cm]current page.north west) {BLOQUE};
  \node[anchor=north west, text=white, font=\bfseries\fontsize{72}{72}\selectfont]
    at ([xshift=0.65cm,yshift=-2.9cm]current page.north west) {#1};
  \node[anchor=north west, text=iubnavy, font=\bfseries\LARGE,
        text width=9.4cm, align=left]
    at ([xshift=6.2cm,yshift=-2.8cm]current page.north west)
    {\hyphenpenalty=10000\exhyphenpenalty=10000 #2};
  \node[anchor=north west, fill=iubnavy, text=iubyellow, rounded corners=1.5mm,
        font=\small\bfseries, inner sep=2mm]
    at ([xshift=6.2cm,yshift=-6.0cm]current page.north west) {Duración: #3};
  \fill[iubnavy] ([xshift=5.2cm]current page.south west) rectangle
        ([yshift=0.35cm]current page.south east);
  \fill[iubyellow] ([xshift=5.2cm,yshift=0.35cm]current page.south west)
        rectangle ([yshift=0.42cm]current page.south east);
  \node[anchor=north east, text=iubnavy, font=\small\bfseries]
    at ([xshift=-0.6cm,yshift=-0.5cm]current page.north east) {IUB · IEL04};
\end{tikzpicture}
\end{frame}}

% =====================================================================
\begin{document}
% =====================================================================

% ---------------------------------------------------------------------
%  PORTADA
% ---------------------------------------------------------------------
\begin{frame}[plain]
\begin{tikzpicture}[remember picture, overlay]
  % Panel institucional
  \fill[iubnavy] (current page.north west) rectangle
        ([xshift=5.6cm]current page.south west);
  \node[anchor=north west, text=iubyellow, font=\bfseries\fontsize{54}{54}\selectfont]
    at ([xshift=0.7cm,yshift=-1.0cm]current page.north west) {IUB};
  \node[anchor=north west, text=white, font=\footnotesize, text width=4.3cm, align=left]
    at ([xshift=0.75cm,yshift=-3.0cm]current page.north west)
    {Institución Universitaria de Barranquilla};
  \draw[iubyellow, line width=1.2pt]
    ([xshift=0.75cm,yshift=-4.25cm]current page.north west) --
    ([xshift=2.6cm,yshift=-4.25cm]current page.north west);
  \node[anchor=north west, text=iubyellow, font=\scriptsize\bfseries,
        text width=4.3cm, align=left]
    at ([xshift=0.75cm,yshift=-4.5cm]current page.north west)
    {Tecnología en Gestión de Sistemas Eléctricos};
  \node[anchor=south west, text=white, font=\scriptsize, text width=4.3cm, align=left]
    at ([xshift=0.75cm,yshift=0.9cm]current page.south west)
    {Ing. Sergio Martinez Castilla\\[1pt]\textcolor{iubyellow}{\tiny smartinezc@unibarranquilla.edu.co}};
  % Bloque de título
  \node[anchor=north west, fill=iubblue, text=white, rounded corners=1.5mm,
        font=\scriptsize\bfseries, inner sep=2mm]
    at ([xshift=6.4cm,yshift=-1.0cm]current page.north west)
    {IEL04 \textperiodcentered{} Semana 13 de 14 \textperiodcentered{} Corte evaluativo 3};
  \node[anchor=north west, text=iubnavy, font=\small, text width=9.0cm, align=left] (a)
    at ([xshift=6.4cm,yshift=-1.9cm]current page.north west)
    {Circuitos Eléctricos I};
  \node[anchor=north west, text=iubnavy, font=\bfseries\fontsize{13}{16}\selectfont,
        text width=9.0cm, align=left] (t)
    at ([yshift=-0.35cm]a.south west)
    {\hyphenpenalty=10000 Cálculo y medición de parámetros en circuitos RLC en serie, en paralelo y mixtos: filtros resonantes, frecuencia de resonancia, ancho de banda y nomenclatura};
  % Franjas de la paleta complementaria
  \fill[iubblue]   ([xshift=6.4cm,yshift=1.1cm]current page.south west)
    rectangle ++(1.6cm,0.18cm);
  \fill[iuborange] ([xshift=8.1cm,yshift=1.1cm]current page.south west)
    rectangle ++(1.6cm,0.18cm);
  \fill[iubgreen]  ([xshift=9.8cm,yshift=1.1cm]current page.south west)
    rectangle ++(1.6cm,0.18cm);
  \node[anchor=north west, text=iubnavy, font=\scriptsize]
    at ([xshift=6.4cm,yshift=0.95cm]current page.south west)
    {Sesión teórico-práctica \textperiodcentered{} 240 min};
\end{tikzpicture}
\end{frame}

% ---------------------------------------------------------------------
\begin{frame}{Punto de partida de la sesión}
\begin{columns}[T,onlytextwidth]
  \column{0.34\textwidth}
\begin{caja}[height=5.9cm]{iubblue}{Continuidad}
\scriptsize \begin{itemize}\setlength{\itemsep}{2pt}
  \item Semana 12: Circuito RLC en paralelo
  \item \textbf{Semana 13: Cálculo y medición de parámetros en circuitos RLC serie y paralelo: unidades, simbología y nomenclatura}
  \item Semana 14: evaluación
\end{itemize}
\end{caja}
  \column{0.34\textwidth}
\begin{caja}[height=5.9cm]{iuborange}{Prerrequisitos}
\scriptsize \begin{itemize}\setlength{\itemsep}{2pt}
  \item Calcular impedancia, corrientes y tensiones de circuitos RLC en serie
  \item Calcular la frecuencia de resonancia, el factor de calidad y el ancho
  \item Convertir una bobina real a su equivalente en paralelo
  \item Aplicar el método general de análisis fasorial (semana 7)
\end{itemize}
\end{caja}
  \column{0.29\textwidth}
\begin{caja}[height=5.9cm]{iubgreen}{Competencia}
\scriptsize \textbf{UC2.} Coordinar actividades asociadas a sistemas eléctricos utilizando fuentes convencionales y no convencionales de energía.
\end{caja}
\end{columns}
\end{frame}

% ---------------------------------------------------------------------
\begin{frame}{Resultados de aprendizaje}
\begin{columns}[c,onlytextwidth]
  \column{0.45\textwidth}
\begin{caja}{iubnavy}{IEL04-RA4}
\footnotesize Calcular un circuito resistivo, inductivo, capacitivo y RLC, sea en serie o paralelo.
\end{caja}
\vspace{1mm}
\begin{cajaplana}{iubblue}
\scriptsize \textbf{CE1.} Identificar los tipos de resistores, inductores y capacitores.\\[2pt]
\textbf{CE2.} Realizar mediciones en un circuito resistivo, inductivo y capacitivo.\\[2pt]
\textbf{CE3.} Realizar mediciones en un circuito RLC.
\end{cajaplana}
  \column{0.52\textwidth}
\centering
\begin{adjustbox}{max width=\linewidth, max totalheight=6.1cm}
\begin{tikzpicture}
  \node[box, fill=iubblue, text width=4.9cm, align=left, font=\tiny, anchor=west] (r0) at (1.25,0.00) {Identificar los componentes de un circuito RLC y los parámetros que caracterizan cada tipo de conexión.};
  \node[tag=iubnavy, font=\tiny\bfseries, minimum width=1.15cm, anchor=east] at (1.1,0.00) {Aplicar};
  \node[box, fill=iuborange, text width=4.9cm, align=left, font=\tiny, anchor=west] (r1) at (1.25,-1.45) {Calcular la respuesta de circuitos RLC mixtos y de filtros resonantes pasabanda y rechazabanda.};
  \node[tag=iubnavy, font=\tiny\bfseries, minimum width=1.15cm, anchor=east] at (1.1,-1.45) {Aplicar};
  \node[box, fill=iubgreen, text width=4.9cm, align=left, font=\tiny, anchor=west] (r2) at (1.25,-2.90) {Diferenciar los parámetros de los circuitos RLC en serie y en paralelo con sus unidades, simbología y nomenclatura.};
  \node[tag=iubnavy, font=\tiny\bfseries, minimum width=1.15cm, anchor=east] at (1.1,-2.90) {Analizar};
  \node[box, fill=iubblue, text width=4.9cm, align=left, font=\tiny, anchor=west] (r3) at (1.25,-4.35) {Medir la frecuencia de resonancia, las frecuencias de corte, el ancho de banda y el factor de calidad de un filtro RLC y evaluar la concordancia con el cálculo.};
  \node[tag=iubnavy, font=\tiny\bfseries, minimum width=1.15cm, anchor=east] at (1.1,-4.35) {Evaluar};
  \draw[flecha, draw=iubnavy!40] (r0.south) -- (r1.north);
  \draw[flecha, draw=iubnavy!40] (r1.south) -- (r2.north);
  \draw[flecha, draw=iubnavy!40] (r2.south) -- (r3.north);
\end{tikzpicture}
\end{adjustbox}
\end{columns}
\end{frame}

% ---------------------------------------------------------------------
\begin{frame}{Ruta de la sesión \textperiodcentered{} 240 minutos}
\centering
\begin{tikzpicture}[x=1cm,y=1cm]
  \node[circle, fill=iubblue, text=white, font=\scriptsize\bfseries, minimum size=0.6cm, inner sep=0] at (0,0.00) {1};
  \node[anchor=west, font=\scriptsize, text width=7.6cm] at (0.45,0.00) {Qué parámetros definen un circuito RLC y cómo se obtienen por cálculo y por medición};
  \fill[iubblue] (8.3,-0.20) rectangle ++(1.84,0.4);
  \node[anchor=west, font=\scriptsize\bfseries] at (10.24,0.00) {20 min};
  \node[circle, fill=iuborange, text=white, font=\scriptsize\bfseries, minimum size=0.6cm, inner sep=0] at (0,-0.76) {2};
  \node[anchor=west, font=\scriptsize, text width=7.6cm] at (0.45,-0.76) {Magnitud común, ley de Kirchhoff, resonancia, Q y comportamiento en frecuencia};
  \fill[iuborange] (8.3,-0.96) rectangle ++(2.76,0.4);
  \node[anchor=west, font=\scriptsize\bfseries] at (11.16,-0.76) {30 min};
  \node[circle, fill=iubgreen, text=white, font=\scriptsize\bfseries, minimum size=0.6cm, inner sep=0] at (0,-1.51) {3};
  \node[anchor=west, font=\scriptsize, text width=7.6cm] at (0.45,-1.51) {Reducción del tanque real y tensión de salida en función de la frecuencia};
  \fill[iubgreen] (8.3,-1.71) rectangle ++(3.22,0.4);
  \node[anchor=west, font=\scriptsize\bfseries] at (11.62,-1.51) {35 min};
  \node[circle, fill=iubblue, text=white, font=\scriptsize\bfseries, minimum size=0.6cm, inner sep=0] at (0,-2.27) {4};
  \node[anchor=west, font=\scriptsize, text width=7.6cm] at (0.45,-2.27) {Frecuencias de corte, ancho de banda y respuesta de los filtros RLC en serie};
  \fill[iubblue] (8.3,-2.47) rectangle ++(3.22,0.4);
  \node[anchor=west, font=\scriptsize\bfseries] at (11.62,-2.27) {35 min};
  \node[circle, fill=iuborange, text=white, font=\scriptsize\bfseries, minimum size=0.6cm, inner sep=0] at (0,-3.03) {5};
  \node[anchor=west, font=\scriptsize, text width=7.6cm] at (0.45,-3.03) {Barrido de frecuencia, criterio de fase, puntos de 70.7 \% y Q = fr/AB};
  \fill[iuborange] (8.3,-3.23) rectangle ++(2.76,0.4);
  \node[anchor=west, font=\scriptsize\bfseries] at (11.16,-3.03) {30 min};
  \node[circle, fill=iubgreen, text=white, font=\scriptsize\bfseries, minimum size=0.6cm, inner sep=0] at (0,-3.79) {6};
  \node[anchor=west, font=\scriptsize, text width=7.6cm] at (0.45,-3.79) {Cuadro final de magnitudes del curso: R, L, C, X, B, Z, Y, fr, Q, AB, P, Q, S y FP};
  \fill[iubgreen] (8.3,-3.99) rectangle ++(2.30,0.4);
  \node[anchor=west, font=\scriptsize\bfseries] at (10.70,-3.79) {25 min};
  \node[circle, fill=iubblue, text=white, font=\scriptsize\bfseries, minimum size=0.6cm, inner sep=0] at (0,-4.54) {7};
  \node[anchor=west, font=\scriptsize, text width=7.6cm] at (0.45,-4.54) {Medir fr, f1, f2, AB y Q del filtro formado con 100 \ensuremath{\Omega}, L3 y C2 y compararlos con el cálculo};
  \fill[iubblue] (8.3,-4.74) rectangle ++(4.60,0.4);
  \node[anchor=west, font=\scriptsize\bfseries] at (13.00,-4.54) {50 min};
  \node[circle, fill=iuborange, text=white, font=\scriptsize\bfseries, minimum size=0.6cm, inner sep=0] at (0,-5.30) {8};
  \node[anchor=west, font=\scriptsize, text width=7.6cm] at (0.45,-5.30) {Síntesis, cierre actitudinal y bibliografía};
  \fill[iuborange] (8.3,-5.50) rectangle ++(1.38,0.4);
  \node[anchor=west, font=\scriptsize\bfseries] at (9.78,-5.30) {15 min};
  \draw[iubnavy, line width=0.6pt] (8.3,0.45) -- (8.3,-5.70);
\end{tikzpicture}
\end{frame}

% =====================================================================
\bloque{01}{Qué parámetros definen un circuito RLC y cómo se obtienen por cálculo y por medición}{20 min}{iubblue}
% =====================================================================

% ---------------------------------------------------------------------
\begin{frame}{{\normalsize Del análisis a la caracterización de un circuito RLC}}
\begin{columns}[c,onlytextwidth]
  \column{0.57\textwidth}
\small
\begin{itemize}\setlength{\itemsep}{4pt}
  \item Al llegar a la última sesión del curso, el estudiante sabe calcular la impedancia de cualquier combinación de resistores, bobinas y condensadores a una frecuencia dada.
  \item Del lado de la medición, reúne los métodos para localizar la resonancia y las frecuencias de corte con el generador y el osciloscopio.
\end{itemize}
  \column{0.4\textwidth}
\begin{cardO}{Nota pedagógica}
\footnotesize Esta sesión es también un repaso para la evaluación final de la semana 14.
\end{cardO}
\end{columns}
\end{frame}

% ---------------------------------------------------------------------
\begin{frame}{Caracterización de un circuito RLC}
\centering
\begin{adjustbox}{max width=\linewidth, max totalheight=6.1cm}
\begin{tikzpicture}
  \node[bG, align=center, font=\scriptsize, text width=3.51cm] (n0) at (1.88,4.53) {Identificar R, L, RW y C};
  \node[bB, align=center, font=\scriptsize, text width=2.36cm] (n1) at (6.66,4.53) {Tipo de conexión};
  \node[bB, align=center, font=\scriptsize, text width=2.80cm] (n2) at (11.09,4.53) {Calcular fr, Q y AB};
  \node[bB, align=center, font=\scriptsize, text width=2.80cm] (n3) at (11.09,1.07) {Barrer la frecuencia};
  \node[bO, align=center, font=\scriptsize, text width=2.58cm] (n4) at (6.66,1.07) {Medir fr, f1 y f2};
  \node[bG, align=center, font=\scriptsize, text width=3.51cm] (n5) at (1.88,1.07) {Comparar y explicar diferencias};
  \draw[flecha] (n0) -- (n1);
  \draw[flecha] (n1) -- (n2);
  \draw[flecha] (n2) -- (n3);
  \draw[flecha] (n3) -- (n4);
  \draw[flecha] (n4) -- (n5);
\end{tikzpicture}
\end{adjustbox}

\vspace{1mm}{\scriptsize Caracterización de un circuito RLC: del cálculo de sus parámetros a su verificación en el laboratorio\par}
\end{frame}

% =====================================================================
\bloque{02}{Magnitud común, ley de Kirchhoff, resonancia, Q y comportamiento en frecuencia}{30 min}{iuborange}
% =====================================================================

% ---------------------------------------------------------------------
\begin{frame}{{\normalsize Comparación sistemática de los circuitos RLC en serie}}
\begin{columns}[c,onlytextwidth]
  \column{0.57\textwidth}
\small
\begin{itemize}\setlength{\itemsep}{4pt}
  \item Esta dualidad no es una curiosidad matemática; es la mejor herramienta para no confundir las fórmulas.
  \item La diferencia más importante está en qué ocurre en la resonancia.
\end{itemize}
  \column{0.4\textwidth}
\begin{cardO}{Nota pedagógica}
\footnotesize Para decidir con rapidez si un circuito es inductivo o capacitivo, conviene preguntar qué elemento «manda» a esa frecuencia: en serie manda el de mayor reactancia; en paralelo, el de menor reactancia, porque es el que conduce más corriente.
\end{cardO}
\end{columns}
\end{frame}

% ---------------------------------------------------------------------
\begin{frame}{Comparación de los circuitos RLC en serie (1/2)}
\centering\scriptsize
\renewcommand{\arraystretch}{1.35}
\rowcolors{2}{iubnavy!6}{white}
\begin{adjustbox}{max width=\linewidth, max totalheight=5.6cm}
\begin{tabularx}{\textwidth}{>{\raggedright\arraybackslash}X>{\raggedright\arraybackslash}X>{\raggedright\arraybackslash}X}
  \rowcolor{iubnavy}
  \hd{Característica} & \hd{RLC en serie (semana 11)} & \hd{RLC en paralelo (semana 12)}\\
  Magnitud común & corriente & tensión\\
  Suma natural & Z = R + j(XL \ensuremath{-} XC) & Y = G + j(BC \ensuremath{-} BL)\\
  Ley de Kirchhoff que reparte & LVK: VS = VR + j(VL \ensuremath{-} VC) & LCK: IT = IR + j(IC \ensuremath{-} IL)\\
  Por debajo de fr & capacitivo & inductivo\\
  Por encima de fr & inductivo & capacitivo\\
  En la resonancia & Z mínima = R; I máxima & Z máxima = Rp; I mínima\\
  Magnitud que se multiplica por Q & tensiones en L y C: Q·VS & corrientes en L y C: Q·IT\\
\end{tabularx}
\end{adjustbox}
\vspace{2mm}
{\scriptsize Comparación de los circuitos RLC en serie y en paralelo\par}
\end{frame}

% ---------------------------------------------------------------------
\begin{frame}{Comparación de los circuitos RLC en serie (2/2)}
\centering\scriptsize
\renewcommand{\arraystretch}{1.35}
\rowcolors{2}{iubnavy!6}{white}
\begin{adjustbox}{max width=\linewidth, max totalheight=5.6cm}
\begin{tabularx}{\textwidth}{>{\raggedright\arraybackslash}X>{\raggedright\arraybackslash}X>{\raggedright\arraybackslash}X}
  \rowcolor{iubnavy}
  \hd{Característica} & \hd{RLC en serie (semana 11)} & \hd{RLC en paralelo (semana 12)}\\
  Factor de calidad & Q = XL / R & Q = R / XL\\
  Una R mayor & reduce Q & aumenta Q\\
  Efecto de RW de la bobina & se suma a R & equivale a RW(Q² + 1) en paralelo\\
  Resonancia medida con el kit & 4.97 kHz & 4.96 kHz\\
  Uso típico & filtro pasabanda en serie, compensación & circuito tanque, rechazo de banda, compensación de FP\\
\end{tabularx}
\end{adjustbox}
\vspace{2mm}
{\scriptsize Comparación de los circuitos RLC en serie y en paralelo\par}
\end{frame}

% =====================================================================
\bloque{03}{Reducción del tanque real y tensión de salida en función de la frecuencia}{35 min}{iubgreen}
% =====================================================================

% ---------------------------------------------------------------------
\begin{frame}{Circuitos RLC mixtos}
\begin{columns}[c,onlytextwidth]
  \column{0.57\textwidth}
\small
\begin{itemize}\setlength{\itemsep}{4pt}
  \item Un circuito RLC mixto muy frecuente es un resistor $R_1$ en serie con un circuito tanque, con la salida tomada en el tanque.
  \item Con $R_1 = 1\ \text{k}\Omega$, la bobina L3 (10.4 mH y 24.6 \ensuremath{\Omega}) y el condensador C2 (97.8 nF), a tres frecuencias.
  \item El circuito es un \textbf{pasabanda} construido con un tanque en paralelo, el dual del pasabanda en serie que se caracteriza en el caso de estudio.
\end{itemize}
  \column{0.4\textwidth}
\begin{caja}{iubblue}{Ecuación clave}
\footnotesize \centering\small \begin{adjustbox}{max width=\linewidth}$\displaystyle \begin{gathered}\mathbf{Z}_p = \frac{1}{j2\pi f C + \dfrac{1}{R_W + j2\pi f L}} \\ \frac{\mathbf{V}_{sal}}{\mathbf{V}_S} = \frac{\mathbf{Z}_p}{R_1 + \mathbf{Z}_p}\end{gathered}$\end{adjustbox}
\end{caja}
\end{columns}
\end{frame}

% ---------------------------------------------------------------------
\begin{frame}{Circuito mixto R1 + tanque real}
\centering\scriptsize
\renewcommand{\arraystretch}{1.35}
\rowcolors{2}{iubnavy!6}{white}
\begin{adjustbox}{max width=\linewidth, max totalheight=5.6cm}
\begin{tabularx}{\textwidth}{>{\raggedright\arraybackslash}X>{\raggedright\arraybackslash}X>{\raggedright\arraybackslash}X>{\raggedright\arraybackslash}X>{\raggedright\arraybackslash}X}
  \rowcolor{iubnavy}
  \hd{f (Hz)} & \hd{Zp (\ensuremath{\Omega})} & \hd{Carácter del tanque} & \hd{ZT (\ensuremath{\Omega})} & \hd{Vsal/VS}\\
  3000 & 308.6 \ensuremath{\angle} 78.8° & inductivo & 1102 & 0.280\\
  4976 (fr) & 4349 \ensuremath{\angle} 0° & resistivo & 5349 & 0.813\\
  8000 & 332.4 \ensuremath{\angle} \ensuremath{-}88.3° & capacitivo & 1063 & 0.313\\
\end{tabularx}
\end{adjustbox}
\vspace{2mm}
{\scriptsize Circuito mixto R1 + tanque real (R1 = 1 k\ensuremath{\Omega}, L = 10.4 mH, RW = 24.6 \ensuremath{\Omega}, C = 97.8 nF), salida en el tanque\par}
\end{frame}

% =====================================================================
\bloque{04}{Frecuencias de corte, ancho de banda y respuesta de los filtros RLC en serie}{35 min}{iubblue}
% =====================================================================

% ---------------------------------------------------------------------
\begin{frame}{Filtros resonantes pasabanda y rechazabanda}
\begin{columns}[c,onlytextwidth]
  \column{0.57\textwidth}
\small
\begin{itemize}\setlength{\itemsep}{4pt}
  \item El ancho de banda se mide entre las dos frecuencias de corte, en las que la salida del pasabanda cae al 70.7 \% de su máximo.
  \item Con $R = 100\ \Omega$, $L = 10\ \text{mH}$ y $C = 0.1\ \mu\text{F}$: $f_r = 5033\ \text{Hz}$ y $R/(4\pi L) = 795.8\ \text{Hz}$, de modo que $f_1 = \sqrt{5033^2 + 795.8^2} - 795.8 \approx 4300\ \text{Hz}$ y $f_2 \approx 5891\ \text{Hz}$.
  \item Las dos curvas son complementarias: donde una tiene su máximo, la otra tiene su mínimo, y ambas se cruzan en las frecuencias de corte, a 0.707.
\end{itemize}
  \column{0.4\textwidth}
\begin{caja}{iubblue}{Ecuación clave}
\footnotesize \centering\small \begin{adjustbox}{max width=\linewidth}$\displaystyle \begin{gathered}f_{1,2} = \sqrt{f_r^{2} + \left(\frac{R}{4\pi L}\right)^{2}} \mp \frac{R}{4\pi L} \\ AB = f_2 - f_1 = \frac{R}{2\pi L} = \frac{f_r}{Q} \\ f_1 f_2 = f_r^{2}\end{gathered}$\end{adjustbox}
\end{caja}
\end{columns}
\end{frame}

% ---------------------------------------------------------------------
\begin{frame}{Respuesta de un RLC en serie}
\begin{columns}[c,onlytextwidth]
  \column{0.62\textwidth}
\centering
\begin{tikzpicture}
  \begin{semilogxaxis}[width=8.4cm, height=5.7cm, xmin=1000, xmax=25000, ymin=0.0093192, ymax=1.0538,
      xlabel={Frecuencia f (Hz)}, ylabel={Vsal / Vent},
      label style={font=\scriptsize, text=iubnavy}, tick label style={font=\tiny, text=iubnavy},
      axis line style={iubnavy}, grid=major, grid style={iubnavy!12},
      legend style={font=\tiny, fill=white}, legend pos=north east]
    \addplot[line width=1.4pt, iubblue] coordinates {(1000,0.065275) (1044,0.068386) (1089.9,0.071668) (1137.9,0.075135) (1188,0.078801) (1240.2,0.082681) (1294.8,0.086794) (1351.8,0.091158) (1411.3,0.095798) (1473.4,0.10074) (1538.2,0.10601) (1605.9,0.11164) (1676.6,0.11767) (1750.3,0.12414) (1827.3,0.13111) (1907.8,0.13863) (1991.7,0.14677) (2079.3,0.15562) (2170.8,0.16527) (2266.4,0.17584) (2366.1,0.18746) (2470.2,0.20032) (2578.9,0.21461) (2692.4,0.2306) (2810.8,0.24861) (2934.5,0.26905) (3063.7,0.29244) (3198.5,0.31945) (3339.2,0.35095) (3486.1,0.38806) (3639.5,0.43226) (3799.7,0.48539) (3966.9,0.54971) (4141.4,0.62749) (4323.7,0.71981) (4513.9,0.8232) (4712.5,0.9232) (4919.9,0.98984) (5136.4,0.99182) (5362.4,0.92805) (5598.4,0.82897) (5844.7,0.72526) (6101.9,0.63218) (6370.4,0.55361) (6650.7,0.4886) (6943.4,0.43492) (7248.9,0.39029) (7567.8,0.35283) (7900.9,0.32105) (8248.5,0.29382) (8611.5,0.27025) (8990.4,0.24966) (9386,0.23153) (9799,0.21544) (10230,0.20106) (10680,0.18813) (11150,0.17644) (11641,0.16582) (12153,0.15613) (12688,0.14724) (13246,0.13906) (13829,0.13151) (14438,0.12451) (15073,0.11801) (15736,0.11196) (16429,0.1063) (17151,0.10102) (17906,0.096061) (18694,0.091405) (19517,0.087026) (20375,0.0829) (21272,0.079008) (22208,0.075331) (23185,0.071853) (24205,0.068561) (25000,0.066206)};
    \addlegendentry{Pasabanda (salida en R)}
    \addplot[line width=1.4pt, iuborange] coordinates {(1000,0.99787) (1044,0.99766) (1089.9,0.99743) (1137.9,0.99717) (1188,0.99689) (1240.2,0.99658) (1294.8,0.99623) (1351.8,0.99584) (1411.3,0.9954) (1473.4,0.99491) (1538.2,0.99437) (1605.9,0.99375) (1676.6,0.99305) (1750.3,0.99226) (1827.3,0.99137) (1907.8,0.99034) (1991.7,0.98917) (2079.3,0.98782) (2170.8,0.98625) (2266.4,0.98442) (2366.1,0.98227) (2470.2,0.97973) (2578.9,0.9767) (2692.4,0.97305) (2810.8,0.9686) (2934.5,0.96313) (3063.7,0.95628) (3198.5,0.9476) (3339.2,0.93639) (3486.1,0.92163) (3639.5,0.90175) (3799.7,0.8743) (3966.9,0.83536) (4141.4,0.77862) (4323.7,0.69417) (4513.9,0.56775) (4712.5,0.38431) (4919.9,0.14219) (5136.4,0.12767) (5362.4,0.37245) (5598.4,0.55929) (5844.7,0.68848) (6101.9,0.77482) (6370.4,0.83278) (6650.7,0.87251) (6943.4,0.90047) (7248.9,0.92069) (7567.8,0.93569) (7900.9,0.94706) (8248.5,0.95586) (8611.5,0.96279) (8990.4,0.96833) (9386,0.97283) (9799,0.97652) (10230,0.97958) (10680,0.98214) (11150,0.98431) (11641,0.98616) (12153,0.98774) (12688,0.9891) (13246,0.99028) (13829,0.99132) (14438,0.99222) (15073,0.99301) (15736,0.99371) (16429,0.99433) (17151,0.99488) (17906,0.99538) (18694,0.99581) (19517,0.99621) (20375,0.99656) (21272,0.99687) (22208,0.99716) (23185,0.99742) (24205,0.99765) (25000,0.99781)};
    \addlegendentry{Rechazabanda (salida en L y C)}
    \addplot[line width=1pt, dashed, iubnavy!70] coordinates {(1000,0.707) (25000,0.707)};
    \addlegendentry{70.7 \%}
    \addplot[line width=1pt, dashed, iubnavy!70] coordinates {(4300,0.0093192) (4300,1.0538)};
    \addlegendentry{f1}
    \addplot[line width=1pt, dashed, iubnavy!70] coordinates {(5891,0.0093192) (5891,1.0538)};
    \addlegendentry{f2}
  \end{semilogxaxis}
\end{tikzpicture}
  \column{0.35\textwidth}
\begin{caja}{iubblue}{Lectura de la gráfica}
\footnotesize Las dos curvas son complementarias: donde una tiene su máximo, la otra tiene su mínimo, y ambas se cruzan en las frecuencias de corte, a 0.707.
\end{caja}
\end{columns}
\end{frame}

% =====================================================================
\bloque{05}{Barrido de frecuencia, criterio de fase, puntos de 70.7 \% y Q = fr/AB}{30 min}{iuborange}
% =====================================================================

% ---------------------------------------------------------------------
\begin{frame}{{\normalsize Métodos para medir la resonancia, el ancho de banda}}
\begin{columns}[c,onlytextwidth]
  \column{0.57\textwidth}
\small
\begin{itemize}\setlength{\itemsep}{4pt}
  \item Caracterizar un circuito resonante en el laboratorio consiste en obtener tres números: la frecuencia de resonancia, el ancho de banda y el factor de calidad.
\end{itemize}
  \column{0.4\textwidth}
\begin{caja}{iubblue}{Ecuación clave}
\footnotesize \centering\small \begin{adjustbox}{max width=\linewidth}$\displaystyle \begin{gathered}AB = f_2 - f_1 \\ Q = \frac{f_r}{AB} \\ f_r \approx \sqrt{f_1 f_2}\end{gathered}$\end{adjustbox}
\end{caja}
\end{columns}
\end{frame}

% ---------------------------------------------------------------------
\begin{frame}{Procedimiento para medir la frecuencia}
\centering
\begin{adjustbox}{max width=\linewidth, max totalheight=6.1cm}
\begin{tikzpicture}
  \node[bG, align=center, font=\scriptsize, text width=3.51cm] (n0) at (1.88,4.53) {Fijar VS en los bornes};
  \node[bB, align=center, font=\scriptsize, text width=2.80cm] (n1) at (6.77,4.53) {Barrer f y hallar el pico};
  \node[bB, align=center, font=\scriptsize, text width=3.03cm] (n2) at (11.43,4.53) {Afinar fr por fase 0°};
  \node[bB, align=center, font=\scriptsize, text width=3.03cm] (n3) at (11.43,1.07) {Buscar 70.7 \% a cada lado};
  \node[bO, align=center, font=\scriptsize, text width=2.13cm] (n4) at (6.77,1.07) {f1 y f2 medidas};
  \node[bG, align=center, font=\scriptsize, text width=3.06cm] (n5) at (1.88,1.07) {AB = f2 \ensuremath{-} f1, Q = fr/AB};
  \draw[flecha] (n0) -- (n1);
  \draw[flecha] (n1) -- (n2);
  \draw[flecha] (n2) -- (n3);
  \draw[flecha] (n3) -- (n4);
  \draw[flecha] (n4) -- (n5);
\end{tikzpicture}
\end{adjustbox}

\vspace{1mm}{\scriptsize Procedimiento para medir la frecuencia de resonancia, las frecuencias de corte, el ancho de banda y el factor de calidad\par}
\end{frame}

% =====================================================================
\bloque{06}{Cuadro final de magnitudes del curso: R, L, C, X, B, Z, Y, fr, Q, AB, P, Q, S y FP}{25 min}{iubgreen}
% =====================================================================

% ---------------------------------------------------------------------
\begin{frame}{Parámetros, unidades, simbología y nomenclatura}
\begin{columns}[c,onlytextwidth]
  \column{0.57\textwidth}
\small
\begin{itemize}\setlength{\itemsep}{4pt}
  \item Esta unidad reúne en un solo cuadro todas las magnitudes que han aparecido desde la semana 1.
  \item Dos coincidencias de nombre deben manejarse con cuidado.
  \item En la \textbf{simbología} de los planos, el resistor lleva el designador R, la bobina L y el condensador C, cada uno seguido de un número correlativo.
\end{itemize}
  \column{0.4\textwidth}
\begin{cardO}{Error frecuente}
\footnotesize Escribir la potencia reactiva en vatios o la aparente en VAR.
\end{cardO}
\end{columns}
\end{frame}

% ---------------------------------------------------------------------
\begin{frame}{{\normalsize Parámetros de los circuitos R, L, C, RC, RL y RLC (1/3)}}
\centering\scriptsize
\renewcommand{\arraystretch}{1.35}
\rowcolors{2}{iubnavy!6}{white}
\begin{adjustbox}{max width=\linewidth, max totalheight=5.6cm}
\begin{tabularx}{\textwidth}{>{\raggedright\arraybackslash}X>{\raggedright\arraybackslash}X>{\raggedright\arraybackslash}X>{\raggedright\arraybackslash}X>{\raggedright\arraybackslash}X}
  \rowcolor{iubnavy}
  \hd{Parámetro} & \hd{Símbolo} & \hd{Unidad} & \hd{Habituales} & \hd{Definición o fórmula}\\
  Resistencia & R & ohmio (\ensuremath{\Omega}) & k\ensuremath{\Omega}, M\ensuremath{\Omega} & V/I en el resistor\\
  Inductancia & L & henrio (H) & µH, mH & v = L di/dt\\
  Resistencia de devanado & RW & ohmio (\ensuremath{\Omega}) & \ensuremath{\Omega} & medida con ohmímetro\\
  Capacitancia & C & faradio (F) & pF, nF, µF & Q/V\\
  Reactancia inductiva & XL & ohmio (\ensuremath{\Omega}) & \ensuremath{\Omega}, k\ensuremath{\Omega} & 2\ensuremath{\pi}fL\\
  Reactancia capacitiva & XC & ohmio (\ensuremath{\Omega}) & \ensuremath{\Omega}, k\ensuremath{\Omega} & 1/(2\ensuremath{\pi}fC)\\
  Impedancia & Z, Z\ensuremath{\angle}\ensuremath{\theta} & ohmio (\ensuremath{\Omega}) & \ensuremath{\Omega}, k\ensuremath{\Omega} & R + j(XL \ensuremath{-} XC) en serie\\
\end{tabularx}
\end{adjustbox}
\vspace{2mm}
{\scriptsize Parámetros de los circuitos R, L, C, RC, RL y RLC: símbolo, unidad, submúltiplos habituales y definición\par}
\end{frame}

% ---------------------------------------------------------------------
\begin{frame}{{\normalsize Parámetros de los circuitos R, L, C, RC, RL y RLC (2/3)}}
\centering\scriptsize
\renewcommand{\arraystretch}{1.35}
\rowcolors{2}{iubnavy!6}{white}
\begin{adjustbox}{max width=\linewidth, max totalheight=5.6cm}
\begin{tabularx}{\textwidth}{>{\raggedright\arraybackslash}X>{\raggedright\arraybackslash}X>{\raggedright\arraybackslash}X>{\raggedright\arraybackslash}X>{\raggedright\arraybackslash}X}
  \rowcolor{iubnavy}
  \hd{Parámetro} & \hd{Símbolo} & \hd{Unidad} & \hd{Habituales} & \hd{Definición o fórmula}\\
  Conductancia & G & siemens (S) & mS, µS & 1/R\\
  Susceptancias & BL, BC & siemens (S) & mS & 1/XL, 1/XC\\
  Admitancia & Y, Y\ensuremath{\angle}\ensuremath{\theta} & siemens (S) & mS & G + j(BC \ensuremath{-} BL) en paralelo\\
  Constante de tiempo & \ensuremath{\tau} & segundo (s) & µs, ms & RC o L/R\\
  Frecuencia de resonancia & fr & hercio (Hz) & kHz, MHz & 1/(2\ensuremath{\pi}\ensuremath{\surd}LC)\\
  Factor de calidad & Q & adimensional & — & XL/R (serie), R/XL (paralelo)\\
  Ancho de banda & AB & hercio (Hz) & kHz & f2 \ensuremath{-} f1 = fr/Q\\
\end{tabularx}
\end{adjustbox}
\vspace{2mm}
{\scriptsize Parámetros de los circuitos R, L, C, RC, RL y RLC: símbolo, unidad, submúltiplos habituales y definición\par}
\end{frame}

% ---------------------------------------------------------------------
\begin{frame}{{\normalsize Parámetros de los circuitos R, L, C, RC, RL y RLC (3/3)}}
\centering\scriptsize
\renewcommand{\arraystretch}{1.35}
\rowcolors{2}{iubnavy!6}{white}
\begin{adjustbox}{max width=\linewidth, max totalheight=5.6cm}
\begin{tabularx}{\textwidth}{>{\raggedright\arraybackslash}X>{\raggedright\arraybackslash}X>{\raggedright\arraybackslash}X>{\raggedright\arraybackslash}X>{\raggedright\arraybackslash}X}
  \rowcolor{iubnavy}
  \hd{Parámetro} & \hd{Símbolo} & \hd{Unidad} & \hd{Habituales} & \hd{Definición o fórmula}\\
  Potencia real & P & vatio (W) & mW, kW & I²R = VI cos \ensuremath{\theta}\\
  Potencia reactiva & Q & voltamperio reactivo (VAR) & mVAR, kVAR & I²X = VI sen \ensuremath{\theta}\\
  Potencia aparente & S & voltamperio (VA) & mVA, kVA & VI\\
  Factor de potencia & FP & adimensional & — & cos \ensuremath{\theta} = P/S\\
\end{tabularx}
\end{adjustbox}
\vspace{2mm}
{\scriptsize Parámetros de los circuitos R, L, C, RC, RL y RLC: símbolo, unidad, submúltiplos habituales y definición\par}
\end{frame}

% =====================================================================
\bloque{07}{Medir fr, f1, f2, AB y Q del filtro formado con 100 \ensuremath{\Omega}, L3 y C2 y compararlos con el cálculo}{50 min}{iubblue}
% =====================================================================

% ---------------------------------------------------------------------
\begin{frame}{{\normalsize Caracterización de un filtro pasabanda RLC en serie}}
\begin{columns}[c,onlytextwidth]
  \column{0.57\textwidth}
\small
\begin{itemize}\setlength{\itemsep}{4pt}
  \item Los tres se verifican de nuevo antes del montaje, con lo que se cubre el criterio CE1.
  \item La resistencia total del lazo es $R = 99.5 + 24.6 = 124.1\ \Omega$.
  \item El caso también muestra la importancia de fijar la tensión en los bornes del circuito.
\end{itemize}
  \column{0.4\textwidth}
\begin{caja}{iubblue}{Ecuación clave}
\footnotesize \centering\small \begin{adjustbox}{max width=\linewidth}$\displaystyle \begin{gathered}f_{1,2}= \sqrt{4990^{2} + 949.6^{2}} \mp 949.6 \\ \approx 5080 \mp 950\ \text{Hz} \ \Rightarrow\ f_1 \\ \approx 4130\ \text{Hz},\ \ f_2 \\ \approx 6029\ \text{Hz}\end{gathered}$\end{adjustbox}
\end{caja}
\end{columns}
\end{frame}

% ---------------------------------------------------------------------
\begin{frame}{{\normalsize Caracterización del filtro pasabanda RLC en serie del kit}}
\centering\scriptsize
\renewcommand{\arraystretch}{1.35}
\rowcolors{2}{iubnavy!6}{white}
\begin{adjustbox}{max width=\linewidth, max totalheight=5.6cm}
\begin{tabularx}{\textwidth}{>{\raggedright\arraybackslash}X>{\raggedright\arraybackslash}X>{\raggedright\arraybackslash}X>{\raggedright\arraybackslash}X>{\raggedright\arraybackslash}X}
  \rowcolor{iubnavy}
  \hd{Parámetro} & \hd{Ideal nominal (R = 100 \ensuremath{\Omega}, sin RW)} & \hd{Cálculo con valores medidos y RW} & \hd{Medido} & \hd{Diferencia con el cálculo}\\
  fr & 5033 Hz & 4990 Hz & 4970 Hz & \ensuremath{-}0.4 \%\\
  Vsal máxima & 2.00 V & 1.60 V & 1.59 V & \ensuremath{-}0.6 \%\\
  f1 & 4300 Hz & 4130 Hz & 4120 Hz & \ensuremath{-}0.2 \%\\
  f2 & 5891 Hz & 6029 Hz & 6050 Hz & +0.3 \%\\
  AB = f2 \ensuremath{-} f1 & 1592 Hz & 1899 Hz & 1930 Hz & +1.6 \%\\
  Q = fr/AB & 3.16 & 2.63 & 2.58 & \ensuremath{-}1.9 \%\\
  \ensuremath{\surd}(f1·f2) & 5033 Hz & 4990 Hz & 4993 Hz & comprobación\\
\end{tabularx}
\end{adjustbox}
\vspace{2mm}
{\scriptsize Caracterización del filtro pasabanda RLC en serie del kit (lecturas de ejemplo)\par}
\end{frame}

% =====================================================================
\bloque{08}{Síntesis, cierre actitudinal y bibliografía}{15 min}{iuborange}
% =====================================================================

% ---------------------------------------------------------------------
\begin{frame}{Síntesis de la sesión}
\begin{columns}[T,onlytextwidth]
  \column{0.32\textwidth}
\begin{cardB}{Comparación sistemática de los circuitos}
\scriptsize Magnitud común, ley de Kirchhoff, resonancia, Q y comportamiento en frecuencia
\end{cardB}
  \column{0.32\textwidth}
\begin{cardO}{Circuitos RLC mixtos}
\scriptsize Reducción del tanque real y tensión de salida en función de la frecuencia
\end{cardO}
  \column{0.32\textwidth}
\begin{cardG}{Filtros resonantes pasabanda}
\scriptsize Frecuencias de corte, ancho de banda y respuesta de los filtros RLC en serie
\end{cardG}
\end{columns}
\vspace{2mm}
\begin{columns}[T,onlytextwidth]
  \column{0.32\textwidth}
\begin{cardB}{Métodos para medir la resonancia, el ancho}
\scriptsize Barrido de frecuencia, criterio de fase, puntos de 70.7 \% y Q = fr/AB
\end{cardB}
  \column{0.32\textwidth}
\begin{cardO}{Parámetros, unidades, simbología}
\scriptsize Cuadro final de magnitudes del curso: R, L, C, X, B, Z, Y, fr, Q, AB, P, Q, S y FP
\end{cardO}
  \column{0.32\textwidth}
\begin{cardG}{Caracterización de un filtro pasabanda RLC}
\scriptsize Medir fr, f1, f2, AB y Q del filtro formado con 100 \ensuremath{\Omega}, L3 y C2 y compararlos con el cálculo
\end{cardG}
\end{columns}
\end{frame}

% ---------------------------------------------------------------------
\begin{frame}{Reflexión para llevar}
\centering
\begin{tikzpicture}
  \node[fill=iubnavy, text=white, rounded corners=6pt, text width=11.5cm, inner sep=12pt, align=center, font=\normalsize] (c) at (0,0) {\itshape «Al cerrar el curso, devolver el kit verificado, rotulado y completo, con un registro de los valores medidos, es la última muestra de responsabilidad con el equipo: el grupo del próximo semestre empezará su trabajo sobre los datos que este grupo deja.»};
  \node[fill=iubyellow, text=iubnavy, rounded corners=3pt, font=\scriptsize\bfseries, inner sep=4pt, anchor=south west] at ([yshift=-4pt]c.north west) {Componente actitudinal};
\end{tikzpicture}
\end{frame}

% ---------------------------------------------------------------------
\begin{frame}{Referencias bibliográficas}
\small
\begin{itemize}\setlength{\itemsep}{8pt}
  \item[{$[1]$}] T. L. Floyd, Principios de circuitos eléctricos, 8.\textordfeminine{} ed. México: Pearson Educación, 2007.
  \item[{$[2]$}] R. L. Boylestad, Introducción al análisis de circuitos, 10.\textordfeminine{} ed. México: Pearson Educación, 2004.
  \item[{$[3]$}] M. F. P. Deorsola y P. Morcelle del Valle, Circuitos eléctricos. Parte 1, 1.\textordfeminine{} ed. La Plata: Editorial de la Universidad de La Plata, 2017.
\end{itemize}
\end{frame}

% ---------------------------------------------------------------------
%  CIERRE
% ---------------------------------------------------------------------
\begin{frame}[plain]
\begin{tikzpicture}[remember picture, overlay]
  \fill[iubnavy] (current page.north west) rectangle (current page.south east);
  \node[anchor=north west, text=iubyellow, font=\bfseries\fontsize{40}{44}\selectfont]
    at ([xshift=1.2cm,yshift=-1.6cm]current page.north west) {\textexclamdown{}Gracias!};
  \node[anchor=north west, text=white, font=\normalsize, text width=14cm, align=left]
    at ([xshift=1.2cm,yshift=-3.4cm]current page.north west)
    {IEL04 \textperiodcentered{} Circuitos Eléctricos I};
  \node[anchor=north west, fill=iubblue, text=white, rounded corners=1.5mm,
        font=\small\bfseries, inner sep=2.2mm, text width=11.5cm]
    at ([xshift=1.2cm,yshift=-4.4cm]current page.north west)
    {Próxima sesión \textperiodcentered{} Semana 14: evaluación};
  \node[anchor=south west, text=iubyellow, font=\small, align=left]
    at ([xshift=1.2cm,yshift=0.9cm]current page.south west)
    {Ing. Sergio Martinez Castilla \quad · \quad smartinezc@unibarranquilla.edu.co};
  \fill[iubblue]   ([xshift=1.2cm,yshift=0.55cm]current page.south west) rectangle ++(1.6cm,0.15cm);
  \fill[iuborange] ([xshift=2.9cm,yshift=0.55cm]current page.south west) rectangle ++(1.6cm,0.15cm);
  \fill[iubgreen]  ([xshift=4.6cm,yshift=0.55cm]current page.south west) rectangle ++(1.6cm,0.15cm);
\end{tikzpicture}
\end{frame}

\end{document}
